\documentclass[11pt]{article}

\usepackage{amsmath, amssymb, amsthm}
\usepackage[margin=1in]{geometry}

\title{Definition and Properties of the Moment Generating Function}
\author{DeepSeek prompted by Riyad Derguini}
\date{}

\begin{document}

\maketitle

\section{Definition: Moment Generating Function}

Let \(X\) be a random variable. The \textbf{Moment Generating Function (MGF)} of \(X\) is a function \(M_X : \mathbb{R} \to [0, \infty]\) defined by:
\[
 M_X(t) = \mathbb{E}[e^{tX}] 
\]
for all \(t \in \mathbb{R}\) for which the expectation is finite. 

If \(X\) is a discrete random variable with probability mass function \(p(x)\), this is:
\[
 M_X(t) = \sum_x e^{tx} p(x) 
\]
If \(X\) is a continuous random variable with probability density function \(f(x)\), this is:
\[
 M_X(t) = \int_{-\infty}^{\infty} e^{tx} f(x) \, dx 
\]

\hrulefill

\section{Why is it called ``Moment Generating''?}

The name comes from the fact that if you differentiate \(M_X(t)\) at \(t = 0\), you generate the \textbf{moments} of the random variable (its expected value, expected square, expected cube, etc.).

We can see this by expanding \(e^{tX}\) using its Maclaurin series:
\[
 e^{tX} = 1 + tX + \frac{(tX)^2}{2!} + \frac{(tX)^3}{3!} + \dots 
\]
Taking the expectation of both sides:
\[
 M_X(t) = \mathbb{E}[e^{tX}] = 1 + t\mathbb{E}[X] + \frac{t^2}{2!}\mathbb{E}[X^2] + \frac{t^3}{3!}\mathbb{E}[X^3] + \dots 
\]

Now, take the \(k\)-th derivative of \(M_X(t)\) and evaluate it at \(t=0\):
\begin{itemize}
    \item \(M_X'(0) = \mathbb{E}[X]\) (the 1st moment, or mean)
    \item \(M_X''(0) = \mathbb{E}[X^2]\) (the 2nd moment)
    \item \(M_X^{(k)}(0) = \mathbb{E}[X^k]\) (the \(k\)-th moment)
\end{itemize}

\hrulefill

\section{Key Properties}

\begin{enumerate}
    \item \textbf{Independence (The most important property for algorithms/ML):}
    If \(X\) and \(Y\) are independent random variables, then the MGF of their sum is the product of their MGFs:
    \[
 M_{X+Y}(t) = \mathbb{E}[e^{t(X+Y)}] = \mathbb{E}[e^{tX}e^{tY}] = \mathbb{E}[e^{tX}]\mathbb{E}[e^{tY}] = M_X(t) M_Y(t) 
\]
    This property allows the multiplication of the MGFs of individual independent elements to obtain the MGF of the whole sum.

    \item \textbf{Uniqueness:}
    If two random variables have the same MGF in an open interval around \(t=0\), they have the exact same probability distribution.
\end{enumerate}

\hrulefill

\section{Connection to the Chernoff Bound}

The MGF is the fundamental tool of the \textbf{Chernoff Method}. The basic inequality states that for any \(t > 0\),
\[
 \Pr(X \ge a) = \Pr(e^{tX} \ge e^{ta}) \le \frac{\mathbb{E}[e^{tX}]}{e^{ta}} = \frac{M_X(t)}{e^{ta}} 
\]
By calculating the MGF, a problem about probabilities is transformed into a problem about algebra. One then chooses \(t\) to minimize the right-hand side, yielding an exponentially decaying bound on the tail probability.

\hrulefill

\section{Examples of MGFs}

\begin{itemize}
    \item \textbf{Rademacher (coin flip, \(\pm 1\) with probability \(1/2\)):}
    \[
 M_X(t) = \frac{1}{2}e^t + \frac{1}{2}e^{-t} = \cosh(t) 
\]
    \item \textbf{Normal Distribution \(\mathcal{N}(\mu, \sigma^2)\):}
    \[
 M_X(t) = \exp\left(\mu t + \frac{1}{2}\sigma^2 t^2\right) 
\]
    \item \textbf{Bernoulli(\(p\)):}
    \[
 M_X(t) = (1-p) + p e^t 
\]
\end{itemize}

\hrulefill

\section{Application to the Discrepancy Proof}

To show that the discrepancy \(D = \sum_{v \in S} \chi(v)\) is concentrated with high probability, one constructs its MGF. 
Because each \(\chi(v)\) is a Rademacher variable, its MGF is \(\cosh(t)\). Because the variables are independent, the MGF of the sum \(D\) is:
\[
 M_D(t) = \prod_{v \in S} \cosh(t) = (\cosh(t))^{|S|} 
\]
Using the inequality \( \cosh(t) \le e^{t^2/2} \), this becomes:
\[
 M_D(t) \le \left(e^{t^2/2}\right)^{|S|} = e^{|S|t^2/2} 
\]
This is the MGF of a Gaussian. By plugging this into the Chernoff bound, one proves that the discrepancy is highly concentrated around \(0\) (with high probability), completing the proof.

\end{document}
